\documentclass[14pt]{extarticle}

\usepackage[a4paper,margin=1in]{geometry}
\usepackage{amsmath,amssymb}
\usepackage{booktabs}
\usepackage{hyperref}

\setlength{\parindent}{0pt}

\def\mathbi#1{\textbf{\em #1}}

\title{Huber Regressor}
\author{Nguyen Phu Nhan}
\date{}

\begin{document}

\maketitle

Prerequisites: Linear Regression. \\

We define and process $\mathbf{X}, \mathbf{y}, \mathbf{w}$ similar to Linear Regression. \\

Huber Regressor extends Linear Regression by replacing the squared loss with the Huber loss, which is less sensitive to outliers. The resulting optimization problem is
\[
\operatorname*{argmin}_{\mathbf{w}}\left(\sum_{i=1}^{n} H_{\epsilon}\left(y_i - \mathbf{x}_i^\top \mathbf{w}\right) + \lambda \|\mathbf{w}\|_2^2\right),
\]

where
\[
H_{\epsilon}(r)=
\begin{cases}
r^2, & |r| \le \epsilon, \\
2\epsilon |r| - \epsilon^2, & |r| > \epsilon.
\end{cases}
\]

Here, $r = y_i - \mathbf{x}_i^\top \mathbf{w}$ is the residual, $\epsilon$ controls the point at which the loss changes from quadratic to linear, and $\lambda \ge 0$ controls the strength of $L_2$ regularization. \\

Unlike ordinary Linear Regression or Ridge, Huber Regressor does not have a closed-form solution for $\mathbf{w}$ because the Huber loss is piecewise-defined. Therefore, $\mathbf{w}$ is usually obtained using iterative optimization methods, such as gradient descent. \\

Let
\[
r_i = y_i - \mathbf{x}_i^\top \mathbf{w}.
\]

The gradient-based update rule is
\[
\mathbf{w}^{(t+1)} = \mathbf{w}^{(t)} + \eta\sum_{i=1}^{n}\psi_{\epsilon}(r_i)\mathbf{x}_i - 2\eta\lambda \mathbf{w}^{(t)},
\]
where
\[
\psi_{\epsilon}(r_i) =
\begin{cases}
2r_i, & |r_i| \le \epsilon, \\
2\epsilon \operatorname{sign}(r_i), & |r_i| > \epsilon.
\end{cases}
\]

Here, $\eta$ is the learning rate, $\epsilon$ controls the threshold between squared loss and absolute loss, and $\lambda$ controls the strength of $L_2$ regularization.
\end{document}